% !TeX root = ../main.tex

\section{Tables}

This section demonstrates the table style used throughout
the technical reports.

The notation and unit conventions introduced in
Section~\ref{sec:equations_units} are also used in the following tables.

\subsection{Design Specifications}

Table~\ref{tab:specifications} summarizes the main design targets.

\begin{table}[htbp]
	\centering
	
	\caption{Example design specifications.}
	\label{tab:specifications}
	
	\begin{tabular}{lll}
		\toprule
		Parameter & Target & Unit \\
		\midrule
		Supply voltage      & 1.8  & V   \\
		DC gain             & 60   & dB  \\
		Unity-gain bandwidth & 100  & MHz \\
		Load capacitance    & 2    & pF  \\
		Bias current        & 2.5  & \textmu A \\
		\bottomrule
	\end{tabular}
	
\end{table}

\subsection{Extended Specification Table}

\begin{table}[htbp]
	\centering
	
	\caption{Extended design specification summary.}
	\label{tab:extended_specifications}
	
	\begin{tabularx}{\textwidth}{l l X}
		\toprule
		Parameter & Target & Description \\
		\midrule
		
		Supply voltage
		& \qty{1.8}{\volt}
		& Nominal supply used during initial transistor-level design. \\
		
		DC gain
		& \qty{60}{\deci\bel}
		& Minimum open-loop low-frequency voltage gain. \\
		
		Unity-gain bandwidth
		& \qty{100}{\mega\hertz}
		& Target small-signal unity-gain frequency under nominal loading. \\
		
		Load capacitance
		& \qty{2}{\pico\farad}
		& Nominal capacitive load connected to the amplifier output. \\
		
		\bottomrule
	\end{tabularx}
	
\end{table}

\subsection{Simulation Results}

\begin{table}[htbp]
	\centering
	
	\caption{Example simulation results.}
	\label{tab:simulation_results}
	
	\begin{tabular}{l S S S}
		\toprule
		Corner
		& {DC Gain (dB)}
		& {GBW (MHz)}
		& {Phase Margin (deg)} \\
		\midrule
		
		TT & 95.21 & 103.4 & 63.2 \\
		SS & 89.43 & 91.7  & 67.8 \\
		FF & 98.82 & 118.6 & 58.9 \\
		
		\bottomrule
	\end{tabular}
	
\end{table}